\newpage
\head{K--1: sorting with tokens}
Each pair uses pancakes 1--4 and the spatula cut from materials page 1; Problems 5--6 use pancakes 1 and 2, brown face as the dark side. Before Problem 1, show the recording once with the printed example $3\,1\,2$, on scrap paper: the top slot is the top pancake, and when the stack is sorted, count the tokens into the flips box and clear them. The adult or partner may write numbers a child says.

Stacks are written from the top; ``flips 3 2'' means flip 3, then flip 2. A bar marks a pancake whose dark side faces up. In the pictures a dark band on a pancake's top edge means its dark side is up. Every count below is the fewest; children's own counts will often be higher, and that is fine.

\sub{Problem 1, page 1: all six stacks of three}
$1\,2\,3$ needs 0; $2\,1\,3$ needs 1 (flip 2); $3\,2\,1$ needs 1 (flip 3); $3\,1\,2$ needs 2 (flips 3 2, the printed example); $2\,3\,1$ needs 2 (flips 2 3); $1\,3\,2$ needs 3 (flips 2 3 2 or 3 2 3). Circle $1\,3\,2$. There are two stacks with each pancake on top, so six in all; the sorted stack counts as one of the six. A child's reason for $1\,3\,2$: ``I tried both first flips, and after either one it still took two more.'' The ``most flips'' depends on fewest counts, so if a box shows more than the number above, ask ``Can you do it with fewer tokens?'' before the circling. Hints: ``How will you know when you have all six?'' ``Is a sorted stack one of the ways?'' Extension: make a stack for your partner that needs exactly two flips.

\sub{Problem 2, page 2: four stacks of four}
$4\,1\,2\,3$: 2 (flips 4 3). $2\,3\,4\,1$: 2 (3 4). $3\,4\,2\,1$: 2 (2 4). $1\,2\,4\,3$: 3, and the only three-flip sort is 4 2 4, which begins by turning the whole stack over. Children often want to swap the 4 and 3 at the bottom, which no flip can do alone. Hint: ``Where must the 4 end up, and how can a flip get it there?''

\sub{Problem 3, page 2: every stack one flip sorts}
$2\,1\,3\,4$, $3\,2\,1\,4$ and $4\,3\,2\,1$, one for each blank. They are complete because the same flip undoes a sorting flip: these are exactly what one flip does to a sorted stack, with the spatula under pancake 2, 3 or 4 (under the top pancake nothing changes). ``Mixed-up'' leaves out $1\,2\,3\,4$. Hint: ``What do you get when you flip a sorted stack once? What if you do the same flip again?'' Extension: find stacks that need exactly two flips (six; table, guide page~\pageref{tab:four}).

\sub{Problem 4, page 3: three hard stacks and a decoy}
$4\,2\,3\,1$: 4 (flips 4 2 3 2). $2\,1\,4\,3$: 3 (3 4 3). $2\,4\,1\,3$: 4 (2 4 3 2). $3\,1\,4\,2$: 4 (2 3 4 2). The three four-flip stacks are the hardest stacks of four; $2\,1\,4\,3$ looks similar and needs only three. Expect first counts of 5 to 8. If a child says a stack ``can't be done in three,'' ask ``How could we be sure?'' and leave it open; the adult can say that grown-ups checked every way. To share: ``Which stack was hardest?''

\sub{Problem 5, page 4: dark sides on two pancakes}
\up1\,2: 1 (flip 1). \up2\,\up1: 1 (flip 2). \up2\,1: 2 (flips 2 1). $2\,1$, both dark down: 3 (1 2 1). The usual slip is flipping a block without turning each pancake over. Hint: ``After that flip, which way is each dark side facing?''

\sub{Problem 6, page 4: the hardest stack of two}
\up1\,\up2, 1 on 2 with both dark sides up, needs 4 (flips 1 2 1 2 or 2 1 2 1); it is the only one. All eight stacks: $1\,2$ needs 0; \up1\,2 and \up2\,\up1 need 1; \up2\,1 and $2$\,\up1 need 2; $2\,1$ and $1$\,\up2 need 3. With only eight stacks (two orders, four ways for the dark sides), the table can build them all. The three blank stacks are working space for candidates, not three answers. Hint: ``How many different stacks of 1 and 2 can you make?''
\newpage
\head{Grades 2--3: fewest flips, a method that always works}
Each pair uses a full set from materials page 1; Problems 6--8 use the brown face as the dark side.
\sub{Problem 1, page 1}
$2\,3\,4\,1$: 2 (flips 3 4). $3\,4\,1\,2$: 3 (2 4 2). $5\,4\,3\,2\,1$: 1 (5). $5\,1\,2\,3\,4$: 2 (5 4). $4\,5\,1\,2\,3$: 3 (2 5 3). $1\,2\,3\,5\,4$: 3 (5 2 5). Each shortest sort is the only one of its length. Hint: ``Where must the biggest pancake end up?''

\sub{Problem 2, page 2: the stack of three that needs the most}
$1\,3\,2$ needs 3 (flips 2 3 2). Why not two: flip 1 does nothing, so the first flip is flip 2, giving $3\,1\,2$, or flip 3, giving $2\,3\,1$. The only stacks one flip from sorted are $2\,1\,3$ and $3\,2\,1$, the sorted stack flipped once, and neither first flip lands on one. With the six stacks laid out, a child can show this by pointing. Breaks (4--5 Problem 3) give only 2 here. Hint: ``What could your first flip be? After it, could one more flip finish?''

\sub{Problem 3, page 2: stacks of six}
$6\,5\,4\,3\,2\,1$: 1 (flip 6). $3\,6\,1\,5\,2\,4$: 6, and only by 5 4 6 3 4 5. $1\,2\,3\,4\,6\,5$: 3 (6 2 6). $5\,2\,6\,3\,1\,4$: 6, and only by 2 3 6 4 2 3. The two scrambled stacks each have six breaks, so six cannot be beaten, but the six-flip sorts are hard to find; the two-flips-per-pancake method uses 8 on each, and 7 to 10 is good work. Their job is to push toward a method that always works. Hint: ``Is there a pancake you can put in its final place and never touch again?''

\sub{Problem 4, page 3: a way for any stack}
Any way that always works, with an honest count. The usual one: flip the biggest pancake not yet in place to the top, then flip it down into place. Children count 10 for six and 18 for ten (two flips for each pancake but the smallest), or 9 and 17 (the last two pancakes need at most one flip). Both are right for their counting, and 9 is the best this way can promise, since $2\,6\,5\,1\,3\,4$ really takes it 9. A child who says 12 and 20 counted flips for the smallest pancake too: ask ``When the others are in place, where is the smallest?'' Why it always works: a placed pancake is below every later spatula. It is not the fastest way; no stack of six needs more than 7 (computer). Hint: ``Can you get the 6 to the bottom from anywhere? Does it ever need to move again?'' Extension: 4--5 Problem 2's ``Explain why it never needs more.''

\sub{Problem 5, page 4: all stacks of four that need four}
$2\,4\,1\,3$, $3\,1\,4\,2$ and $4\,2\,3\,1$. To be sure, work outward from $1\,2\,3\,4$ (table, guide page~\pageref{tab:four}): a stack made from $1\,2\,3\,4$ with three flips can be sorted with three, by doing those flips in reverse order. One flip makes 3 stacks, a second makes 6 new ones and a third 11 more; with $1\,2\,3\,4$ that is 21 of the 24, and the three never made are the answers. Pool it: each child follows one first flip, and each new stack goes on a sticky note in one of three columns, 1, 2 and 3 flips, so a repeat is easy to spot. Crossing off stacks that someone sorts in three finds the same three candidates, but only the outward list shows that they cannot be done. Hint: ``If you mix up a sorted stack with three flips, can you always sort it again in three?'' Extension: how many stacks need exactly 1, 2 and 3 flips? (3, 6, 11.)

\sub{Problem 6, page 5: dark sides on two pancakes}
As K--1 Problem 5: \up1\,2 needs 1 (flip 1); \up2\,\up1 needs 1 (flip 2); \up2\,1 needs 2 (flips 2 1); $2\,1$ needs 3 (1 2 1).

\sub{Problem 7, page 5: the hardest stack of two}
\up1\,\up2, 4 flips (1 2 1 2), the only one; the three blanks are working space for one answer. The full list of eight is under K--1 Problem 6.

\sub{Problem 8, page 6: stacks of three with dark sides}
\up3\,1\,2: 2 (flips 3 2). $3\,1\,2$, all dark down: 3 (1 3 2). $3\,2\,1$, all dark down: 5, only by 2 3 2 3 2. \up1\,\up2\,\up3: 6 (2 3 2 3 2 3 or 3 2 3 2 3 2), one of the two hardest stacks of three. Expect long sorts on the last two and accept any honest count. Hint: ``Which pancake can you put in its place, dark side down, first?''
\newpage
\head{Grades 4--5: bounds, breaks and certificates}
The pair and the referee use full sets from materials page 1; Problems 7--9 use the brown face as the dark side.
\sub{Problem 1, page 1}
$5\,1\,2\,3\,4$: 2 (flips 5 4). $4\,5\,1\,2\,3$: 3 (2 5 3). $1\,2\,3\,5\,4$: 3 (5 2 5). $3\,2\,1\,5\,4$: 3 (4 5 4). $2\,1\,4\,3\,6\,5$: 5, only by 3 5 6 5 3. The last looks easy (three swapped pairs) and is the hardest on the page; after Problem 3 it is the example that has 3 breaks and needs 5.

\sub{Problem 2, page 2: a way for any stack, and why}
The two-flips-per-pancake way needs at most $2(n-2)+1$: 9 for six, 37 for twenty. Why never more: each placed pancake is below every later spatula, so it never moves; placing a pancake takes at most two flips (none to bring it up if it is already on top, none at all if it is already in place); and when only 1 and 2 remain they need at most one flip. A child who counts 10 and 38 (or 12 and 40) has a correct, looser bound. The way really uses 9 on some stacks ($2\,6\,5\,1\,3\,4$), but it is not the fastest. Hint: ``After the biggest pancake is at the bottom, what problem is left?'' Extension: every stack of five can be sorted in five (computer), so place the 6 with two flips and finish the top five in five: 7 for six, which is exactly the worst case for six. The same idea gives $2\cdot15+5=35$ for twenty.

\sub{Problem 3, page 3: breaks}
$3\,5\,2\,1\,4$ needs 4, only by flips 2 5 4 3: $3\,5\,2\,1\,4 \to 5\,3\,2\,1\,4 \to 4\,1\,2\,3\,5 \to 3\,2\,1\,4\,5 \to 1\,2\,3\,4\,5$, with 4, 3, 2, 1, 0 breaks. The explanation: the lifted pancakes keep their neighbours (in reverse order) and nothing below the spatula changes, so the only new touching pair is at the spatula, and one flip can remove at most one break. Four breaks need four flips; the plate matters because a whole-stack flip changes the bottom pair. Common slips: counting 3 on 2 as a break, or forgetting the plate. Hint: ``When you flip, which pancakes get a new neighbour?'' Extension: does every stack need exactly as many flips as it has breaks? ($1\,3\,2$ and $4\,2\,3\,1$ say no.)

\sub{Problem 4, page 4: the stacks of four that need four}
$2\,4\,1\,3$, $3\,1\,4\,2$ and $4\,2\,3\,1$. Each cannot be done in three: $2\,4\,1\,3$ and $3\,1\,4\,2$ have four breaks (the only stacks of four that do). $4\,2\,3\,1$ has three, so a three-flip sort would have to remove a break with every flip. Only flip 4 does, giving $1\,3\,2\,4$ with two breaks, and from there no flip removes one: $3\,1\,2\,4$ and $2\,3\,1\,4$ keep two, and flip 4 goes back. That all are found needs the other 21 stacks shown to take three or fewer: the outward layers (guide page~\pageref{tab:four}) kept in sticky-note columns as in 2--3 Problem 5, or a pooled set of sorts. Common slip: keeping plate 6 from Problem 3's example; for four pancakes the plate is 5, so a 4 at the bottom is not a break. Hint: ``For $4\,2\,3\,1$, which first flips remove a break?''

\sub{Problem 5, page 5: a worst stack of five}
Five flips is the most for five pancakes, and 20 stacks need five. Twelve of them have five breaks, so for those, breaks prove ``not fewer'' and the child's five-flip sort completes the answer: $1\,3\,5\,2\,4$, $1\,4\,2\,5\,3$, $2\,4\,1\,5\,3$, $2\,5\,3\,1\,4$, $3\,1\,5\,2\,4$, $3\,5\,1\,4\,2$, $3\,5\,2\,4\,1$, $4\,1\,3\,5\,2$, $4\,2\,5\,1\,3$, $4\,2\,5\,3\,1$, $5\,2\,4\,1\,3$, $5\,3\,1\,4\,2$. Every five-break stack needs exactly five. Eight more need five with fewer breaks: $1\,3\,2\,5\,4$ (three breaks) and $1\,4\,3\,5\,2$, $1\,5\,3\,2\,4$, $1\,5\,3\,4\,2$, $2\,4\,3\,5\,1$, $2\,5\,1\,4\,3$, $3\,1\,5\,4\,2$, $5\,1\,3\,2\,4$ (four). For $1\,5\,3\,2\,4$ every first flip leaves four breaks, so it needs at least $1+4$; the others need a deeper tree. That no stack of five needs six is a computer fact; Problem 2's way promises only seven. Hint: ``How many breaks can a stack of five have? Can you build one with that many?''

\needspace{13\baselineskip}
\sub{Problem 6, page 5: $4\,6\,2\,5\,1\,3$ in six?}
No; it needs 7. With six breaks, a six-flip sort would remove a break with every flip. Following only break-removing flips (breaks in brackets), every branch gets stuck:
\begin{center}\small\begin{tabular}{l}
$4\,6\,2\,5\,1\,3$ (6)\\
\quad flip 3: $2\,6\,4\,5\,1\,3$ (5)\\
\qquad flip 4: $5\,4\,6\,2\,1\,3$ (4); flip 2: $4\,5\,6\,2\,1\,3$ (3); flip 5: $1\,2\,6\,5\,4\,3$ (2), stuck\\
\qquad flip 5: $1\,5\,4\,6\,2\,3$ (4); flip 4: $6\,4\,5\,1\,2\,3$ (3); flip 6: $3\,2\,1\,5\,4\,6$ (2), stuck\\
\quad flip 5: $1\,5\,2\,6\,4\,3$ (5); flip 2: $5\,1\,2\,6\,4\,3$ (4)\\
\qquad flip 3: $2\,1\,5\,6\,4\,3$ (3), stuck\\
\qquad flip 4: $6\,2\,1\,5\,4\,3$ (3); flip 3: $1\,2\,6\,5\,4\,3$ (2), stuck; or flip 6: $3\,4\,5\,1\,2\,6$ (2), stuck
\end{tabular}\end{center}
Hint: ``If six flips must remove six breaks, what must each flip do?'' A seven-flip sort: 2 3 4 2 6 3 2. With $5\,3\,6\,1\,4\,2$ it is one of the two stacks of six that need seven (computer). An adult can record the tree while the pair finds the break-removing flips.

\sub{Problem 7, page 6: stacks of three with dark sides}
\up3\,1\,2: 2 (flips 3 2). $1\,2$\,\up3: 3 (3 1 3). $3\,1\,2$: 3 (1 3 2). \up1\,\up2\,$3$: 4 (1 2 1 2). $3\,2\,1$: 5 (2 3 2 3 2).

\sub{Problem 8, page 6: the hardest stack of three}
Six flips; exactly two stacks need six, so at most two of the three blanks hold answers: \up1\,\up2\,\up3 (2 3 2 3 2 3) and $1$\,\up2\,\up3 (2 1 2 3 1 3). Problem 7's $3\,2\,1$ (five) is the natural first candidate. That these need six and nothing needs seven is a computer fact: the 48 stacks of three need 0 to 6 flips, 1, 3, 6, 12, 18, 6 and 2 of them. Hint: ``Which stack in Problem 7 was hardest? Can you make one harder?'' Extension: with four pancakes, $1\,2\,3\,4$ with every dark side up needs eight, the most for four (computer).

\sub{Problem 9, page 7: a way with dark sides}
Bring the biggest pancake not in place to the top, turn it over alone if its dark side faces down, then flip it down; it lands dark side down. That is at most three flips per pancake and one for the last: $3n-2=16$ for six. The true worst case for six is 12 (computer), so cleverer ways exist. Hint: ``When the biggest pancake is on top, which way must its dark side face before you flip it down?''

\head{Sources and evidence}
Harry Dweighter (Jacob E.\ Goodman), Problem E2569, \emph{American Mathematical Monthly} 82 (1975), p.\ 1010, posed the problem. William H.\ Gates and Christos H.\ Papadimitriou, ``Bounds for sorting by prefix reversal,'' \emph{Discrete Mathematics} 27 (1979), pp.\ 47--57, prove the $(5n+5)/3$ bound. David S.\ Cohen and Manuel Blum, ``On the problem of sorting burnt pancakes,'' \emph{Discrete Applied Mathematics} 61 (1995), study dark sides. OEIS A058986, ``Sorting by prefix reversal,'' lists the worst-case values. None of these was consulted for this guide (no local copies; the OEIS could not be reached); the claims attributed to them come from the theme outline. Every count, list and sort printed here was recomputed by an exhaustive search of every stack (plain up to 7 pancakes, dark sides up to 6) in \texttt{check\_guide.py}, against the final pages. Library links: Week 3's return visit sorts by neighbour swaps, Week 23 sorts with fixed comparator bars, and Week 3 runs permutations as machines; this theme adds prefix reversals and a counted lower bound. Record actual preparation time, card or foam and whether children could see a stack's order, whether they keep the whole-block rule, how long their sorts take, and where an adult had to decide.
